\section{Refinement of histories and conservation of unity}
\label{nuclear:1}
\subsection{Generation of emissions and realization domain}

The construction starts from APP emissions and their enriched prolongations. An emission preserves the two sheets and their divisions

\[
\delta^\diamond=\operatorname{res}^\diamond+9\operatorname{quo}^\diamond,
\qquad \diamond\in\{\Sigma,\Pi\}.
\]

TRIT produces orientation and carry through

\[
\Delta(\epsilon)=\operatorname{or}(\epsilon)+3\kappa(\epsilon),
\qquad \operatorname{or}(\epsilon)\in\{-1,0,+1\}.
\]

TPK transports the complete pair of sheets and updates memory through

\[
U_{\epsilon_2\epsilon_1}=U_{\epsilon_2}U_{\epsilon_1},\qquad
\mathsf U_\epsilon(m)=\mathsf C_\epsilon m+\kappa(\epsilon),\qquad
\kappa(\epsilon_2\epsilon_1)
=\kappa(\epsilon_2)+\mathsf C_{\epsilon_2}\kappa(\epsilon_1).
\]

The cocycle makes affine composition associative. The path boundary telescopes, but its ordered record is an ordered word, not the sum of its entries. The odometer increments vertical memory when the phase returns; return does not identify the complete states. The reduced catalogue transfer does not replace the action of emissions on the tree; the construction in \cref{iv:cont:historias,iv:cont:concatenacion} preserves this distinction.

The enriched state preserves the intermediate cells, both sheets, remainders, quotients, orientation, carry, clock, memory, boundary, path, and ordered record. An admissible emission jointly updates these fields. Truncation erases exactly its final contribution. A neutral prolongation gives a section of truncation but records a step: it is not an absence of history. This yields finiteness, nonemptiness, and surjectivity at every level, followed by total survival.

This development stops at the analytic realization of this system and its memory. No target numerical values or conventional constants enter as inputs. The following probabilistic coefficients arise solely from the number of effective outgoing emissions.

\subsection{Refinement on the complete multisection}

Fix an actual enriched prefix \(x\) of depth \(k\). Write

\[
H_n=\operatorname{Term}_{k,k+n}(x),\qquad
\rho_n:H_{n+1}\longrightarrow H_n,
\qquad \Omega_x=\varprojlim_n(H_n,\rho_n).
\]

Each element of \(H_n\) is a complete prolongation to that horizon, with its emission witnesses. For \(h\in H_n\), the successors are \(h\epsilon\), for all emissions \(\epsilon\in\operatorname{Out}(h)\). Operation multiplicities are preserved even when reduced terminal data coincide. By \cref{iv:cont:medida},

\[
p_h(\epsilon)=\frac1{d(h)}>0,\qquad
d(h)=\#\operatorname{Out}(h),\qquad
\mu_{n+1}(h\epsilon)=\mu_n(h)p_h(\epsilon).
\]

Therefore,

\begin{equation}
\sum_{\epsilon\in\operatorname{Out}(h)}\mu_{n+1}(h\epsilon)=\mu_n(h),
\qquad (\rho_n)_*\mu_{n+1}=\mu_n.
\label{nuc01:eq:1}
\end{equation}

Proposition~\ref{iv:cont:medida} constructs the unique probability measure \(\mu_x\) on \(\Omega_x\) with these cylinder masses; every nonempty cylinder has positive mass. The measure neither selects a history nor imposes a filter subsequent to the existence of the five constructions.

\begin{proposition}[Refinement of observables is isometric and unital]
\label{nuc01:pullback}

Let \(\mathscr L_n=L^2(H_n,\mu_n)\) and

\[
R_n:\mathscr L_n\longrightarrow\mathscr L_{n+1},\qquad
(R_nf)(h\epsilon)=f(h).
\]

Then \(R_n^*R_n=I\), \(R_n1=1\), and its adjoint is

\begin{equation}
(E_ng)(h)=\sum_{\epsilon\in\operatorname{Out}(h)}p_h(\epsilon)g(h\epsilon).
\label{nuc01:eq:2}
\end{equation}

In particular, \(E_nR_n=I\), \(E_n1=1\), and \(R_nE_n\) is the orthogonal projection onto the observables depending only on the prefix.
\end{proposition}

\begin{proof}
Applying \eqref{nuc01:eq:1},

\[
\|R_nf\|_{n+1}^2
=\sum_h\mu_n(h)|f(h)|^2\sum_\epsilon p_h(\epsilon)
=\|f\|_n^2.
\]

The same grouping in the inner product gives \eqref{nuc01:eq:2}. The statements concerning unity hold pointwise; the projection statement follows from \(E_n=R_n^*\) and \(E_nR_n=I\). Moreover, \(R_n(fg)=R_nf\,R_ng\) and \(R_n\bar f=\overline{R_nf}\): on finite observables, it is also a unital algebra homomorphism.
\end{proof}

Upon iteration, \(R_{m,n}=R_{m-1}\cdots R_n\) coincides with pullback along the composed truncation. The cylindrical pullbacks \(R_{\infty,n}\) realize the inductive limit of these Hilbert spaces in \(L^2(\Omega_x,\mu_x)\). Their union is dense: cylinders form an algebra generating the sigma-algebra, and simple functions from that algebra are dense in \(L^2\). Constructing the limit requires no choice of a trajectory.

\subsubsection{The same isometry in the history basis}

The unitary identification \(W_nf(h)=\sqrt{\mu_n(h)}f(h)\) maps \(\mathscr L_n\) to \(\ell^2(H_n)\). Conjugating \(R_n\) gives

\begin{equation}
\widehat R_n\delta_h
=\sum_{\epsilon\in\operatorname{Out}(h)}\sqrt{p_h(\epsilon)}\,\delta_{h\epsilon}.
\label{nuc01:eq:3}
\end{equation}

The vectors on the right have disjoint supports for distinct predecessors, because \(\rho_n(h\epsilon)=h\). Each column has norm one by normalization. This directly proves isometry on the enriched histories already generated, without imposing that a numerical projection preserve all memory. The unit vector transforms as

\[
\Omega_n=\sum_{h\in H_n}\sqrt{\mu_n(h)}\,\delta_h,
\qquad \widehat R_n\Omega_n=\Omega_{n+1}.
\]

No artificial copy of \(h\) is appended to a result in order to recover it: \(h\epsilon\) is precisely the prolongation already produced by the APP–TRIT–TPK domain.

\subsection{Correlative prolongation of the five operations}

The construction in \cref{iv:cont:mapa-correlativo} defines a single record \(\mathfrak e_\epsilon\) for each emission. Its spectral, solenoidal, gauge, cohomological, and determinantal operations depend on this same emission, with the same source, target, orientation, carry, boundary, path, and ordered record. The cellular action satisfies

\[
\partial\operatorname{Cell}(\epsilon)=\operatorname{Cell}(\epsilon)\partial,
\qquad
\operatorname{Cell}(\epsilon_2\epsilon_1)
=\operatorname{Cell}(\epsilon_2)\operatorname{Cell}(\epsilon_1).
\]

The complex and its determinant line belong to the same object; the constructed filling satisfies \(\partial h_*=z_--z_+\). Totality of the constructor does not require a specialized trivialization to be acyclic or invertible. Joint existence is thus preserved before any specialized reader.

The naturality in \cref{iv:rectangular-natural,iv:cont:memoria-natural} is expressed jointly as

\begin{equation}
u_{(\ell,N'),(k,N)}\mathfrak D_{\ell,N'}(\widehat x_\ell)
=\mathfrak D_{k,N}(\tau_{\ell,k}\widehat x_\ell).
\label{nuc01:eq:4}
\end{equation}

Its proof does not identify an isolated fine edge with a fictitious coarse edge: it restricts the emission history and preserves the actions of its entries. Transport composes over the same cylinder probabilities; balances use conditional expectation and its corrector; incidence and ports use their cellular maps; homology and determinant use the same map of complexes. These are not five limits of separately selected histories.

Two directions must be distinguished. State and coefficient maps run from fine to coarse. The pullbacks \(R_n\) of Proposition~\ref{nuc01:pullback} run from coarse observables to fine observables. Expectation \(E_n\) returns in the reduction direction. This distinction makes it possible to realize \eqref{nuc01:eq:4} analytically without reversing its domains.

The terminal construction in \cref{iv:cont:terminal,iv:cont:ensamblaje,iv:cont:determinantes} prolongs the complete terminals and simultaneously preserves polar telescoping, Gram and losses, incidence with measure, the return cone, and Smith–Bockstein data. Promoting one of these actions to a bounded operator requires the corresponding norm and bound (see the passage to the limit proved at the end of this section). A compatible finite collection does not become a bounded operator on the limit merely by being named as such.

\subsection{Memory fibres: Gram criterion and conservation of unity}

For a nonscalar realization, assign to each prefix \(h\) a Hilbert fibre \(\mathscr M_h\), and to its emission a typed linear transport
\(I_\epsilon:\mathscr M_h\to\mathscr M_{h\epsilon}\). This linear operator is not identified with the affine update \(\mathsf U_\epsilon\) without specifying the representation realizing it.

On \(\mathscr H_n=\bigoplus_{h\in H_n}\mathscr M_h\), define

\[
(V_n\xi)_{h\epsilon}=\sqrt{p_h(\epsilon)}I_\epsilon\xi_h.
\]

Disjointness of the successors proves the exact identity of \cref{iv:cont:gram}:

\begin{equation}
\left.V_n^*V_n\right|_{\mathscr M_h}
=\sum_{\epsilon\in\operatorname{Out}(h)}p_h(\epsilon)I_\epsilon^*I_\epsilon.
\label{nuc01:eq:5}
\end{equation}

The necessary and sufficient criterion for \(V_n\) to be isometric is that the sum in \eqref{nuc01:eq:5} equal \(I\) at every predecessor. Each \(I_\epsilon\) being isometric is sufficient. It is also necessary if all \(I_\epsilon\) are already contractions: each \(I-I_\epsilon^*I_\epsilon\) is then positive, and a sum with strictly positive weights can vanish only when every summand vanishes. The Gram criterion characterizes isometry of the total map; isometry of each transport is a sufficient condition, which is also necessary within the class of contractions.

For the scalar histories of \eqref{nuc01:eq:3}, \(I_\epsilon=1\), so the condition has been proved. In additional analytic fibres, \eqref{nuc01:eq:5} is retained and the precise assumption allowing passage to isometry is identified.

If the fibres have distinguished unit vectors \(u_h\) and \(I_\epsilon u_h=u_{h\epsilon}\), then

\begin{equation}
V_n\Bigl(\bigoplus_h\sqrt{\mu_n(h)}u_h\Bigr)
=\bigoplus_{h\epsilon}\sqrt{\mu_{n+1}(h\epsilon)}u_{h\epsilon}.
\label{nuc01:eq:6}
\end{equation}

Isometry alone does not assert \eqref{nuc01:eq:6}: preserving the norm and preserving a distinguished unit vector are distinct properties.

\subsection{Symmetries and readers on the common domain}

A coherent reversible symmetry of the tree transports all emissions and their multiplicities. Invariance of multiplicities in \cref{iv:cont:medida} proves that it preserves degrees and cylinder probabilities. It may map \(\Omega_x\) to \(\Omega_{gx}\); only when it fixes the prefix does it act within the same prolongation fibre.

Therefore,

\[
(K_gf)(g\omega)=f(\omega)
\]

defines a unitary map \(L^2(\Omega_x,\mu_x)\to L^2(\Omega_{gx},\mu_{gx})\). Naturality of truncation implies \(K_{g,n+1}R_n=R'_nK_{g,n}\), and the corresponding identity for expectations follows by taking adjoints. In memory fibres, one additionally requires the square

\[
S_{g,h\epsilon}I_\epsilon=I_{g\epsilon}S_{g,h}
\]

with \(S_{g,h}\) unitary. This is the intertwiner that turns a tree symmetry into a symmetry of this memory representation; it is not attributed merely by name to every subsequent exceptional group.

For the double readout, take the initial gauge and its prefix transport specified in \cref{exc:doble-lectura,exc:limite-inverso}. Work on a common multisection \(\Omega_x\) where the readers are defined; for the Archimedean reader, this includes the nested compact intervals and diameter contraction specified in \eqref{eq:cilindro}. This is a premise concerning the reader realization, rather than a new existence criterion for the joint constructor. The finite identities are

\begin{equation}
r_{n,m}Q_n=Q_m\rho_{n,m},\qquad
a_{n,m}R^{\rm arq}_n=R^{\rm arq}_m\rho_{n,m}.
\label{nuc01:eq:7}
\end{equation}

They determine the joint readout at the limit. If \(L=(\operatorname{ev}_{\rm arq},Q_\infty)\) and \(\nu=L_*\mu_x\), then

\[
C_L:L^2(Y,\nu)\longrightarrow L^2(\Omega_x,\mu_x),
\qquad C_Lf=f\circ L,
\]

is an isometry, since equality of norms is the definition of the pushforward measure. It preserves unity, and its ranges consist of the observables measurable with respect to the information read out by \(L\). For compatible finite readers, \eqref{nuc01:eq:7} also implies consistency of the pushforward measures.

This isometry does not imply that the reader distinguishes every state. The construction in \cref{exc:limite-inverso} distinguishes the fact that the same visible stream with different hidden memory may have the same incidence. The adjoint of \(C_L\) is expectation with respect to the reader; \(C_LC_L^*\) preserves exactly this observable part. If an enlarged realization \((L,M)\) is injective and Borel, with a measurable inverse on its image—for example, a continuous injection from the compact history space into an appropriate Hausdorff space—its pullback does identify the whole of \(L^2\). The required memory is the memory already preserved in the state, not a new retrospective selection.

\subsection{Passage-to-the-limit theorem for the observable/memory charge}

This section specifies the types in the connection with a memory conservation identity. Canonical refinement of observables supplies the required isometric maps. An additional realization in memory fibres uses the Gram criterion of \eqref{nuc01:eq:5}.

Let \((\mathscr H_n,R_n)\) and \((\mathscr K_n,S_n)\) be two inductive systems of Hilbert spaces with isometric inclusions. For each \(n\), let

\[
J_n:\mathscr H_n\to\mathscr K_n\text{ be isometric},\quad
P_n=J_nJ_n^*,\quad U_n\text{ be unitary on }\mathscr K_n,
\]

and let \(\widetilde A_n\) be bounded and self-adjoint, with \(\sup_n\|\widetilde A_n\|<\infty\). Define

\[
A_n=J_n^*\widetilde A_nJ_n,\quad
T_n=J_n^*U_nJ_n,\quad
\eta_n=(I-P_n)U_nJ_n.
\]

Assume conservation and separation at every level,

\begin{equation}
U_n^*\widetilde A_nU_n=\widetilde A_n,
\qquad [\widetilde A_n,P_n]=0,
\label{nuc01:eq:8}
\end{equation}

and the following refinement squares:

\begin{equation}
J_{n+1}R_n=S_nJ_n,\qquad
P_{n+1}S_n=S_nP_n,
\label{nuc01:eq:9}
\end{equation}

\begin{equation}
U_{n+1}S_n=S_nU_n,\qquad
\widetilde A_{n+1}S_n=S_n\widetilde A_n.
\label{nuc01:eq:10}
\end{equation}

\begin{theorem}
These collections induce bounded operators on the inductive limits, with \(J_\infty\) isometric and \(U_\infty\) unitary, and preserve

\begin{equation}
A_\infty
=T_\infty^*A_\infty T_\infty
 +\eta_\infty^*\widetilde A_\infty\eta_\infty.
\label{nuc01:eq:11}
\end{equation}

The same identity holds at every finite depth. If \(\widetilde A_n\geq0\), the retained charge is positive. Without positivity, a self-adjoint form is preserved, rather than necessarily a positive energy.
\end{theorem}

\begin{proof}
Decompose \(U_nJ_n=J_nT_n+\eta_n\). The two terms belong to \(P_n\mathscr K_n\) and \((I-P_n)\mathscr K_n\). Commutation in \eqref{nuc01:eq:8} makes the cross terms of the form \(\widetilde A_n\) vanish. Applying the other identity in \eqref{nuc01:eq:8},

\[
A_n=J_n^*U_n^*\widetilde A_nU_nJ_n
=T_n^*A_nT_n+\eta_n^*\widetilde A_n\eta_n.
\]

The second equality in \eqref{nuc01:eq:9}, together with the first, implies
\(J_{n+1}^*S_n=R_nJ_n^*\): write
\(P_{n+1}S_n=S_nJ_nJ_n^*=J_{n+1}R_nJ_n^*\) and multiply by \(J_{n+1}^*\).
By \eqref{nuc01:eq:9}–\eqref{nuc01:eq:10},

\[
T_{n+1}R_n=R_nT_n,\qquad
\eta_{n+1}R_n=S_n\eta_n,\qquad
A_{n+1}R_n=R_nA_n.
\]

The operators are thus defined on the algebraic unions of the systems. Isometries and the uniform bound allow extension by continuity to the completions. Compatibility of \(U_n^*\) also follows from compatibility of \(U_n\), because both levels are unitary; the limit is therefore unitary, not merely isometric. The projections \(P_n\) induce a projection whose range is the closure of the union of the ranges of \(J_n\); it is \(J_\infty J_\infty^*\). All identities persist first on finite-level vectors and then by density. This proves \eqref{nuc01:eq:11}.
\end{proof}

Compatibility of \(J_n\) alone does not replace the second equality in \eqref{nuc01:eq:9}. A refinement preserving the observable inclusion while mixing its complement with it may alter the memory defect.

\subsubsection{Explicit compatible nonadic realization}

Suppose \(C_n\) is unitary on \(\mathscr H_n\), natural with respect to \(R_n\), and \(A_n\) is bounded and self-adjoint, uniformly bounded, natural, and satisfies \([A_n,C_n]=0\). On \(\mathscr K_n=\mathscr H_n^{\oplus9}\), define

\[
J_nv=\tfrac13(v,\ldots,v),\quad
S_n=R_n^{\oplus9},\quad
U_n(v_1,\ldots,v_9)=(C_nv_9,v_1,\ldots,v_8),
\quad \widetilde A_n=I_9\otimes A_n.
\]

The factor \(1/3\) normalizes the nine components. The formulas directly verify \eqref{nuc01:eq:8}–\eqref{nuc01:eq:10}, including the projection. The compression is \(T_n=(8I+C_n)/9\), and

\begin{equation}
\eta_n^*\widetilde A_n\eta_n
=A_n-T_n^*A_nT_n
=\frac8{81}(I-C_n)^*A_n(I-C_n).
\label{nuc01:eq:12}
\end{equation}

The final equality uses \(C_n^*A_nC_n=A_n\). By the theorem, \eqref{nuc01:eq:12} also holds at the limit. Nonadic refinement or complete history refinement transports the charge and its memory when \(C_n\) and \(A_n\) satisfy the indicated squares.

For \(A=I\), finite telescoping gives

\[
\|v\|^2=\|T^Nv\|^2+\sum_{j=0}^{N-1}\|\eta T^jv\|^2.
\]

Thus, \(v\mapsto(\eta v,\eta Tv,\ldots,\eta T^{N-1}v,T^Nv)\) is an isometry. The corresponding charge formula replaces the norms by the forms of \(A\) and \(\widetilde A\). The terminal term is not removed without proving convergence to zero; depth compatibility alone does not prove this temporal stability.

If \(D_{N,n}\) denotes this isometry at depth \(n\) and temporal horizon \(N\), the preceding squares also prove

\[
(S_n^{\oplus N}\oplus R_n)D_{N,n}
=D_{N,n+1}R_n.
\]

Indeed, each component satisfies \(\eta_{n+1}T_{n+1}^jR_n=S_n\eta_nT_n^j\), and the terminal component satisfies \(T_{n+1}^NR_n=R_nT_n^N\). Finite-horizon memory accumulation and resolution refinement commute. The same calculation proves equivariance for a symmetry represented by unitaries intertwining \(J,U,P,\widetilde A\); it does not extend the action to symmetries for which these squares have not been constructed.
